% sections/16_architecture_overview.tex
\section{System Architecture Overview}\label{sec:architecture-overview}

The \ISI system is organised into four layers.  Each layer encapsulates a
distinct set of responsibilities and invariants.  The layers are strictly
ordered: each layer depends only on the layer below it; no upward or
circular dependency exists.

% -------------------------------------------------------------------------
\subsection*{Layer~1\enspace Mathematical Layer}

\begin{description}[style=sameline, leftmargin=1em]
  \item[Scope.]  Concentration measurement, aggregation, and composite
    construction.
  \item[Core operations.]  Herfindahl--Hirschman concentration
    (\cref{eq:hhi}), cross-channel aggregation
    (\cref{eq:dual-channel,eq:weighted-b,eq:fuel-avg}), arithmetic
    composite mean (\cref{eq:composite}), four-tier classification
    (\cref{tab:thresholds}).
  \item[Invariants.]  All formulas are fixed for the lifetime of a major
    version.  Axis weights are equal.  Rounding precision is eight decimal
    places, round-half-to-even.  No stochastic or discretionary element
    enters the computation.
  \item[State coverage.]  S0\,--\,S4 in the formal state model
    (\cref{sec:state-model}).
\end{description}

% -------------------------------------------------------------------------
\subsection*{Layer~2\enspace Deterministic Computation Layer}

\begin{description}[style=sameline, leftmargin=1em]
  \item[Scope.]  Platform-independent determinism, canonical serialisation,
    and prohibited-operation enforcement.
  \item[Core operations.]  Rounding gate (\textbf{T5}), canonical
    UTF-8/LF/lexicographic-JSON serialisation
    (\cref{sec:det:serialisation}), country-code ordering, locale
    independence, prohibited-operation enforcement
    (\cref{sec:det:prohibited}).
  \item[Invariants.]  D-1 through D-10
    (\cref{sec:det:invariants}).  IEEE~754 binary64 arithmetic
    (\cref{sec:det:floating-point}).  Hash-breakage clause: any
    serialisation deviation changes the snapshot hash.
  \item[State coverage.]  Transition T5 (rounding gate) at boundary
    S4\,\(\to\)\,S5.
\end{description}

% -------------------------------------------------------------------------
\subsection*{Layer~3\enspace Integrity Layer}

\begin{description}[style=sameline, leftmargin=1em]
  \item[Scope.]  Cryptographic binding, tamper detection, and
    non-repudiation.
  \item[Core operations.]  Country-level SHA-256 hashing (\textbf{T6}),
    snapshot hash construction (\textbf{T7}), Ed25519 signing
    (\textbf{T8}), manifest binding (\textbf{C-3}), hash cascade
    (\cref{sec:snapshot:cascade}).
  \item[Invariants.]  Collision resistance and preimage resistance of
    SHA-256.  Existential unforgeability of Ed25519.  Immutable Snapshot
    Law (\cref{sec:snapshot:immutability}).  Tamper propagation
    (\textbf{C-5}).
  \item[State coverage.]  S5\,--\,S8 (post-rounding through signed
    snapshot).
  \item[Security model.]  Adversary classes A-1 through A-6, security
    goals G-1 through G-4 (\cref{sec:threat-model}).
\end{description}

% -------------------------------------------------------------------------
\subsection*{Layer~4\enspace Governance Layer}

\begin{description}[style=sameline, leftmargin=1em]
  \item[Scope.]  Versioning law, registry management, immutability policy,
    key lifecycle, and publication protocol.
  \item[Core operations.]  Three-part version numbering
    (\cref{sec:gov:versioning}), append-only registry (\textbf{TA-2}),
    change protocol (\cref{sec:gov:change}), key rotation
    (\cref{sec:threat:key-rotation}), archival policy
    (\cref{sec:gov:archival}).
  \item[Invariants.]  Version-definition immutability.  No retroactive
    recomputation.  Published snapshots are permanent.  Superseded keys
    remain accessible.
  \item[State coverage.]  Post-S8 policy constraints and S9 (serving).
\end{description}

\bigskip

The four layers collectively ensure that the \ISI system is
mathematically specified (Layer~1), deterministically reproducible
(Layer~2), cryptographically verifiable (Layer~3), and institutionally
governed (Layer~4).  No layer may be removed or weakened without a
major-version increment.
